\documentclass[10pt,a4paper]{amsart}
\usepackage[margin=19mm]{geometry}
\usepackage{amsmath,amssymb,mathtools}
\usepackage[scr=rsfs,cal=euler]{mathalfa}
\usepackage{xfrac}
\usepackage[unicode,hidelinks,bookmarks=false]{hyperref}
\newcommand{\ie}{i.e.\ }
\newcommand{\cf}{cf.\ }
\newcommand{\resp}{respectively\ }
\newcommand{\Subsection}{\subsection}
\providecommand{\nobreakdash}{\nobreak\hbox{-}}
\setlength{\emergencystretch}{2em}
\multlinegap0pt
\usepackage{amssymb}
\newtheorem{prop}{Proposition}
\newcommand{\Ker}{{\rm \,Ker}}
\newcommand{\dn}{{\rm \,dn}}
\newcommand{\sn}{{\rm \,sn}}
\title{On the orbital stability of periodic snoidal waves for the $\phi^4-$equation}
\author{B.S. Lonardoni, F. Natali}
\date{Source: arXiv:2506.05547v4 (CC BY 4.0)}
\begin{document}
\maketitle

\noindent\emph{Source and licence.} On the orbital stability of periodic snoidal waves for the $\phi^4-$equation. Version arXiv:2506.05547v4 (CC BY 4.0), distributed by arXiv under the Creative Commons Attribution 4.0 International licence, http://creativecommons.org/licenses/by/4.0/, as recorded in the arXiv licence metadata for this exact version and shown on the abstract page https://arxiv.org/abs/2506.05547v4. Full licence notice: \texttt{licenses/arxiv-2506.05547-CC-BY-4.0.txt}.\par\smallskip

\noindent\textit{Editorial scope.} Counted unit: the article's prop leman1 (source0.tex lines 679--681). The article's complete original proof of it is delivered with it (source0.tex lines 682--783, one \texttt{proof} environment of 102 lines). No mathematics is written, repaired or re-derived: the statement and the proof are the article's own lines, in the article's own notation, and no retained line is substituted or re-flowed.  Retained uncounted as the source-local context the proof needs: lines 196--204; lines 655--657; lines 659--661.  Nothing was omitted from any retained span, so the package carries no omission record.  The retained lines cite 1 works, whose entries are transcribed verbatim from the article's own bibliography source in the e-print and delivered with short editorial labels recorded in the item's reference mapping.  No retained line contains a figure environment, an image-inclusion command or drawing code, so no image is delivered and the package declares no asset dependency.  Editorial formatting only, in addition: an AMS \texttt{amsart} preamble that restates the article's own declarations and macros used by the retained lines, namely the environments \texttt{prop} on separate counters, together with the standard packages those lines need.  The retained environments print in this document as Proposition 1 in order of appearance, whereas the article prints the same items with the numbering of its own full document.  The complete native source of the article as distributed in the arXiv e-print for this exact version is bundled unchanged as \texttt{sources/179405/source0.tex}.
\begin{equation}\displaystyle\label{opconstrained2}\mathcal{L}_{\Pi}= \mathcal{G}''(h,ch')-\displaystyle \left(
\displaystyle	\begin{array}{ccc}
		\frac{3}{L}\int_{0}^{L} h^2 \cdot \; dx& & 0 \\\\
		0  & & 0
	\end{array}\right)=\displaystyle\mathcal{L}-\displaystyle \left(
\begin{array}{ccc}
\frac{3}{L}\int_{0}^{L} h^2\cdot \; dx & & 0 \\\\
0  & & 0
\end{array}\right),\end{equation} where $\mathcal{L}$ is given by

 \begin{equation}\label{indexformula12}
 	\text{n}(\mathcal{L}_{{1\Pi}})=\text{n}(\mathcal{L}_1)-{\rm n}_0-{\rm z}_0
 \end{equation}

 \begin{equation}\label{indexformula123}
 	\text{z}(\mathcal{L}_{{1\Pi}})=\text{z}(\mathcal{L}_1)+{\rm z}_0,
 \end{equation}

\begin{prop}\label{leman1}
 Let $L\in (0,2\pi)$ be fixed. The linear operator $\mathcal{L}_{\Pi}$ in $(\ref{opconstrained2})$ has no negative eigenvalues and $(h',ch'')$ is a simple eigenfunction associated with the zero eigenvalue. 
\end{prop}

\begin{proof}
	
	
In order to count the negative eigenvalues, it is necessary to note that $\mathcal{L}_{\Pi}$ is the constrained operator $\mathcal{L}$ defined in $\mathbb{L}_{per,m}^2$ with constrained space $$S= [(1,0),(0,1)] \subset \Ker(\mathcal{L})^{\perp}=[(h',ch'')]^\perp,$$ 
such that $\mathcal{L}_{\Pi}\big |_{S^\perp}=\mathcal{L}$. On the other hand, corresponding to the constrained set $S$, we define the matrix
 
 \begin{equation}\label{matrixD}
 	D=\left[\begin{array}{llll}(\mathcal{L}^{-1}(1,0),(1,0))_{\mathbb{L}^2_{per}}& & (\mathcal{L}^{-1}(1,0),(0,1))_{\mathbb{L}^2_{per}}\\\\
 		(\mathcal{L}^{-1}(1,0),(0,1))_{\mathbb{L}^2_{per}}& & (\mathcal{L}^{-1}(0,1),(0,1))_{\mathbb{L}^2_{per}}\end{array}\right].
 	\end{equation}
Since $\mathcal{L}(0,1)=(0,1)$ and $(1,1)=(1,0)+(0,1)$, we have $$\mathcal{L}^{-1}(1,0)=\mathcal{L}^{-1}(1,1)-\mathcal{L}^{-1}(0,1)=\mathcal{L}^{-1}(1,1)-(0,1)$$ and
 \begin{equation}\label{D}
\left( \mathcal{L}^{-1}(1,0), (1,0) \right)_{\mathbb{L}^2_{per}}%= (\tilde{f},1)_{L^2_{per}}  
= ( \mathcal{L}_1^{-1} 1,1)_{L^2_{per}}.%=:D_1.
\end{equation}
Furthermore,
\begin{equation}\label{D12}
		\left( \mathcal{L}^{-1}(1,0), (0,1) \right)_{\mathbb{L}^2_{per}}=\left(\mathcal{L}^{-1}(1,1),(0,1)\right)_{\mathbb{L}^2_{per}}-\left(\mathcal{L}^{-1}(0,1),(0,1)
		\right)_{\mathbb{L}^2_{per}}=L-L=0\end{equation}
and 
\begin{equation}\label{D123}
	\left( \mathcal{L}^{-1}(0,1), (0,1) \right)_{\mathbb{L}^2_{per}}=\left((0,1),(0,1)\right)_{\mathbb{L}^2_{per}}=L.\end{equation}
According to \eqref{D}, \eqref{D12} and \eqref{D123}, it follows that $D$ can be expressed as $$D=\left[\begin{array}{cc}D_1&  0\\
	0&  L\end{array}\right],$$ where $D_1=(\mathcal{L}_{1}^{-1}1,1)_{L_{per}^2}$. The computation of $D_1$ requires finding, since $\Ker(\mathcal{L}_1)=[h']$, an element $\tilde{f}\in H_{per}^2$ satisfying $\mathcal{L}_1\tilde{f}=1$. 
	
	To obtain an appropriate periodic function $\tilde{f}$, we take the following steps: let us consider the first and the fifth eigenvalues of $\mathcal{L}_1$, and their corresponding eigenfunctions, given respectively by 
	\begin{align*}
		\lambda_0=\frac{1+k^2-2\sqrt{1-k^2+k^4}}{1+k^2}, \quad f_0(x)=1-[1+k^2-\sqrt{1-k^2+k^4}]\sn^2(bx;k),
	\end{align*}
	and 
	\begin{align*}
		\lambda_4=\frac{1+k^2+2\sqrt{1-k^2+k^4}}{1+k^2}, \quad f_4(x)=1-[1+k^2+\sqrt{1-k^2+k^4}]\sn^2(bx;k),
	\end{align*}where $ b =\frac{1}{\sqrt{\omega(1+k^2)}}=\frac{4K(k)}{L}$.
	Under these conditions, defining $B_1=(1+k^2+\sqrt{1-k^2+k^4})$ and $B_2 =-(1+k^2-\sqrt{1-k^2+k^4})$, we have that
	\begin{align}\label{eq1.autovalor}
		B_1f_0=&B_1-[(1+k^2)^2-(1-k^2+k^4)]\sn^2(bx;k)= B_1-3k^2\sn^2(bx;k)
	\end{align}
	and
	\begin{align}\label{eq2.autovalor}
		B_2f_4=B_2+[(1+k^2)^2-(1-k^2+k^4)]\sn^2(bx;k)=B_2+3k^2\sn^2(bx;k).
	\end{align}
It follows from \eqref{eq1.autovalor} and \eqref{eq2.autovalor} that 
\begin{align}\label{B1B2}
B_1f_0+B_2f_4=2\sqrt{1-k^2+k^4}. 
\end{align}
\indent Therefore, by the definition of $\lambda_0$ and $\lambda_4$, and using the equality $(\ref{B1B2})$, we deduce 
\begin{equation}\label{eq3.autovalor}
\mathcal{L}_1(\lambda_4 B_1f_0+\lambda_0B_2f_4)= 2\lambda_0\lambda_4\sqrt{1-k^2+k^4}.
\end{equation}
 Since $\lambda_0\lambda_4\sqrt{1-k^2+k^4}$ is nonzero for $k\in (0,1)$, we obtain by \eqref{eq3.autovalor} that $\mathcal{L}_1\tilde{f}=1$, or equivalently, $\tilde{f}=\mathcal{L}_1^{-1}1$, where $\tilde{f}\in  H_{per}^2$ is defined by $$\tilde{f}=\frac{ 1 }{2\lambda_0\lambda_4\sqrt{1-k^2+k^4}}(	\lambda_4 B_1f_0+\lambda_0B_2f_4).$$
Hence, \begin{align}\label{eq4.autovalor}
D_1=(\mathcal{L}_{1}^{-1}1,1)_{L_{per}^2}=(\tilde{f},1)_{L_{per}^2}=\frac{\lambda_4 B_1(f_0,1)_{L_{per}^2}+\lambda_0B_2(f_4,1)_{L_{per}^2}}{2\lambda_0\lambda_4\sqrt{1-k^2+k^4} }.
\end{align}
On the other hand, using the periodicity of the even function $\displaystyle\dn(u+2K(k);k)= \dn(u;k)$ and \cite[Formula 110.07]{byrd}, we obtain the relation  
$\displaystyle
	 	(\sn^2(bx;k),1)_{L_{per}^2}
	 	=\frac{4(K(k)-E(k))}{bk^2},
$
 where $E(k)=\displaystyle\int_0^{\frac{\pi}{2}}\sqrt{1-k^2\sin(\theta)^2}d\theta$ is the complete elliptic integral of the second kind.% defined as $E(k)=\displaystyle\int_0^{\frac{\pi}{2}}\sqrt{1-k^2\sin(\theta)^2}d\theta$.
 \\ Thus, it follows that
\begin{align}
	\label{eq5.autovalor}
	(f_0,1)_{L_{per}^2}=L-\frac{4B_2}{b}\frac{(E(k)-K(k))}{k^2},
\end{align}
and 
\begin{align}
	\label{eq6.autovalor}
	(f_4,1)_{L_{per}^2}=L+\frac{4B_1}{b}\frac{(E(k)-K(k))}{k^2}.
\end{align}
 Therefore, by \eqref{eq5.autovalor}, \eqref{eq6.autovalor}, and the fact that $b=\frac{4K(k)}{L}$, we obtain
  \begin{equation}\label{eq8.autovalor}
 	\lambda_4B_1	(f_0,1)_{L_{per}^2}+\lambda_0 B_2	(f_4,1)_{L_{per}^2}= 6L\sqrt{1-k^2+k^4}+ \frac{12L \sqrt{1-k^2+k^4}}{(1+k^2)}\frac{(E(k)-K(k))}{ K(k)}.
 \end{equation} 
\indent Next, since $2\lambda_0\lambda_4=\frac{-6(1-k^2)^2}{(1+k^2)^2}$, we conclude from \eqref{eq4.autovalor} and \eqref{eq8.autovalor} that
 \begin{equation}\label{eq9.autovalor}\begin{array}{lllll}
 \displaystyle	D_1=(\mathcal{L}_{1}^{-1}1,1)_{L_{per}^2}&=&\displaystyle \frac{6L}{2\lambda_0\lambda_4} + \frac{12L  }{(1+k^2)}\frac{(E(k)-K(k))}{ K(k)2\lambda_0\lambda_4 }\\\\&=&\displaystyle -\frac{L(1+k^2)^2}{(1-k^2)^2} - \frac{2L (1+k^2) }{(1-k^2)^2}\frac{(E(k)-K(k))}{K(k)}\\\\&=&\displaystyle\frac{-L(1+k^2)}{(1-k^2)^2}\bigg\{ (1+k^2)+2 \frac{(E(k)-K(k))}{ K(k)} \biggl \}.\end{array}
 \end{equation}
 By \cite[Formula 710.00 and 710.02]{byrd}, it follows that
  \begin{equation}\label{eq10.autovalor}
 	  (1-k^2)K(k)<E(k)<K(k), \ \mbox{for all}\ k\in (0,1).
 \end{equation}
 Using \eqref{eq9.autovalor} and \eqref{eq10.autovalor} we obtain $	D_1=(\mathcal{L}_{1}^{-1}1,1)_{L_{per}^2}<0$. 
 
% \begin{align}
% 	\lambda_4A+ \lambda_0B=6\sqrt{1-k^2+k^4}
% \end{align}
% 
%  \begin{align}
% 	AB(\lambda_0-\lambda_4)= \frac{12k^2\sqrt{1-k^2+k^4}}{1+k^2}
% \end{align}
  
 
	
\indent Therefore, using $(\ref{indexformula12})$ and $(\ref{indexformula123})$, it follows that
\begin{equation*}
		{\rm{n}}(\mathcal{L}_{1\Pi}) = 	{\rm{n}}(\mathcal{L}_{1})-{\rm{n}}_0-{\rm{z}}_0=1 - 1 - 0 = 0 \quad \mbox{and}\quad  	{\rm{z}}(\mathcal{L}_{1\Pi})={\rm{z}}(\mathcal{L}_{1\ })+{\rm{z}}_0=1+0=1.
\end{equation*}
Associated with the full linear operator $\mathcal{L}_{\Pi}$, we have demonstrated that
\begin{equation*}
	 {\rm{n}}(\mathcal{L}_{\Pi})=	{\rm{n}}(\mathcal{L})-{\rm{n}}_0-{\rm{z}}_0=1-1-0=0  \quad \mbox{and}\quad  {\rm{z}}(\mathcal{L}_{\Pi})={\rm{z}}(\mathcal{L})+{\rm{z}}_0=1+0=1,
\end{equation*}  as requested.
\end{proof}

\begin{thebibliography}{99}
\bibitem[byrd]{byrd} P.F. Byrd and M.D. Friedman, \textit{Handbook of Elliptic Integrals for Engineers and Scientists}, Second edition, Springer-Verlag, New York, 1971.
\end{thebibliography}
\end{document}
